% ============================================================================
%  preambule.tex — configuration typographique, couleurs, boîtes, diagrammes
%  Compilation : XeLaTeX (ou LuaLaTeX) + Biber
% ============================================================================

% ---------------------------------------------------------------- géométrie
\usepackage[a4paper,top=2.3cm,bottom=2.2cm,left=2.4cm,right=2.4cm,
            headheight=14pt,footskip=1.0cm]{geometry}

% ---------------------------------------------------------------- polices
\usepackage{fontspec}
\usepackage{unicode-math}
\defaultfontfeatures{Ligatures=TeX}

% Chargement défensif : la police de secours est celle que toute installation
% TeX possède, ce qui évite qu'un rapport cesse de compiler ailleurs.
\IfFontExistsTF{TeX Gyre Pagella}
  {\setmainfont{TeX Gyre Pagella}}
  {\setmainfont{Latin Modern Roman}}
\IfFontExistsTF{TeX Gyre Heros}
  {\setsansfont{TeX Gyre Heros}[Scale=MatchLowercase]}
  {\setsansfont{Latin Modern Sans}[Scale=MatchLowercase]}
\IfFontExistsTF{DejaVu Sans Mono}
  {\setmonofont{DejaVu Sans Mono}[Scale=0.84]}
  {\setmonofont{Latin Modern Mono}[Scale=MatchLowercase]}
\IfFontExistsTF{TeX Gyre Pagella Math}
  {\setmathfont{TeX Gyre Pagella Math}}
  {\setmathfont{Latin Modern Math}}

\usepackage[french]{babel}
\frenchsetup{SmallCapsFigTabCaptions=false}
\usepackage{csquotes}
\usepackage{microtype}

% ---------------------------------------------------------------- couleurs
% Reprises de lib/core/constants/app_colors.dart : le rapport porte l'identité
% visuelle du produit qu'il décrit.
\usepackage{xcolor}
\definecolor{emcblue}{HTML}{1E3A8A}   % primaryBlue
\definecolor{emcorange}{HTML}{C2410C} % primaryOrange (AA sur blanc)
\definecolor{emcred}{HTML}{B91C1C}    % dangerRedStrong
\definecolor{emcslate}{HTML}{0F172A}  % textPrimary
\definecolor{emcgrey}{HTML}{64748B}   % textSecondary
\definecolor{emcborder}{HTML}{E2E8F0}
\definecolor{emcbg}{HTML}{F8FAFC}
\definecolor{emcgreen}{HTML}{15803D}

% ---------------------------------------------------------------- tableaux
\usepackage{booktabs}
\usepackage{tabularx}
\usepackage{multirow}
\usepackage{array}
\usepackage{longtable}
\newcolumntype{L}[1]{>{\raggedright\arraybackslash}p{#1}}
\newcolumntype{C}[1]{>{\centering\arraybackslash}p{#1}}
\renewcommand{\arraystretch}{1.18}
\usepackage[table]{colortbl}

% ---------------------------------------------------------------- listes
\usepackage{enumitem}
\setlist[itemize]{leftmargin=1.15em,itemsep=1pt,topsep=2pt,parsep=0pt}
\setlist[enumerate]{leftmargin=1.5em,itemsep=1pt,topsep=2pt,parsep=0pt}

% ---------------------------------------------------------------- figures
\usepackage{graphicx}
\usepackage{float}
% Les flottants remplissent les pages ; sans cela, un [H] trop haut laisse
% une demi-page blanche derrière lui.
\renewcommand{\topfraction}{0.92}
\renewcommand{\bottomfraction}{0.7}
\renewcommand{\textfraction}{0.06}
\renewcommand{\floatpagefraction}{0.75}
\setcounter{topnumber}{3}
\setcounter{bottomnumber}{2}
\setcounter{totalnumber}{5}
% Espacement des flottants resserré : le rapport tient en quinze pages.
\setlength{\intextsep}{7pt plus 2pt minus 2pt}
\setlength{\textfloatsep}{8pt plus 2pt minus 2pt}
\setlength{\floatsep}{7pt plus 2pt minus 2pt}
\setlength{\abovecaptionskip}{4pt}
\setlength{\belowcaptionskip}{0pt}
\usepackage[font=small,labelfont={bf,color=emcblue},
            justification=centering,skip=6pt]{caption}
\usepackage{subcaption}
\captionsetup[sub]{font=footnotesize,labelfont=bf}

% ---------------------------------------------------------------- TikZ
\usepackage{tikz}
\usetikzlibrary{shapes.geometric,shapes.multipart,shapes.symbols,arrows.meta,
                positioning,calc,fit,backgrounds,matrix,chains,
                decorations.pathreplacing,decorations.markings,patterns,quotes}
\usepackage{pgfgantt}

% ---------------------------------------------------------------- boîtes
\usepackage[most]{tcolorbox}
\tcbuselibrary{listings,breakable,skins}

\newtcolorbox{decision}[1][]{%
  enhanced,breakable,colback=emcblue!4,colframe=emcblue,
  boxrule=0pt,leftrule=2.5pt,arc=2pt,left=6pt,right=6pt,top=5pt,bottom=5pt,
  fonttitle=\sffamily\bfseries\small,coltitle=emcblue,
  attach boxed title to top left={yshift=-2mm,xshift=6mm},
  boxed title style={colback=white,boxrule=0pt,arc=1pt},
  title={Décision technique},#1}

\newtcolorbox{vigilance}[1][]{%
  enhanced,breakable,colback=emcorange!5,colframe=emcorange,
  boxrule=0pt,leftrule=2.5pt,arc=2pt,left=6pt,right=6pt,top=5pt,bottom=5pt,
  fonttitle=\sffamily\bfseries\small,coltitle=emcorange,
  attach boxed title to top left={yshift=-2mm,xshift=6mm},
  boxed title style={colback=white,boxrule=0pt,arc=1pt},
  title={Point de vigilance},#1}

\newtcolorbox{constat}[1][]{%
  enhanced,breakable,colback=emcbg,colframe=emcborder,
  boxrule=0.8pt,arc=3pt,left=7pt,right=7pt,top=5pt,bottom=5pt,#1}

% ---------------------------------------------------------------- code
\usepackage{listings}
\lstdefinelanguage{Dart}{
  morekeywords={abstract,as,assert,async,await,break,case,catch,class,const,
    continue,default,deferred,do,dynamic,else,enum,export,extends,external,
    factory,false,final,finally,for,get,if,implements,import,in,is,late,
    library,mixin,new,null,on,operator,part,required,rethrow,return,sealed,
    set,static,super,switch,sync,this,throw,true,try,typedef,var,void,while,
    with,yield,String,int,double,bool,List,Map,Future,Widget,record},
  sensitive=true,
  morecomment=[l]{//},morecomment=[s]{/*}{*/},morecomment=[l]{///},
  morestring=[b]',morestring=[b]"}
\lstdefinelanguage{TypeScript}{
  morekeywords={abstract,any,as,async,await,boolean,break,case,catch,class,
    const,constructor,continue,declare,default,delete,do,else,enum,export,
    extends,false,finally,for,from,function,get,if,implements,import,in,
    instanceof,interface,let,new,null,number,of,private,protected,public,
    readonly,return,set,static,string,super,switch,this,throw,true,try,type,
    typeof,undefined,var,void,while,yield},
  sensitive=true,
  morecomment=[l]{//},morecomment=[s]{/*}{*/},
  morestring=[b]',morestring=[b]",morestring=[b]`}

\lstdefinestyle{emc}{
  basicstyle=\ttfamily\scriptsize\color{emcslate},
  keywordstyle=\color{emcblue}\bfseries,
  commentstyle=\color{emcgrey}\itshape,
  stringstyle=\color{emcgreen},
  numbers=left,numberstyle=\ttfamily\tiny\color{emcgrey},numbersep=8pt,
  showstringspaces=false,breaklines=true,breakatwhitespace=true,
  tabsize=2,columns=fullflexible,keepspaces=true,
  frame=leftline,framerule=1pt,rulecolor=\color{emcborder},
  xleftmargin=14pt,framexleftmargin=14pt,
  aboveskip=8pt,belowskip=6pt,
  literate={é}{{\'e}}1 {è}{{\`e}}1 {ê}{{\^e}}1 {à}{{\`a}}1 {ç}{{\c c}}1
           {É}{{\'E}}1 {ù}{{\`u}}1 {û}{{\^u}}1 {î}{{\^i}}1 {ô}{{\^o}}1
           {«}{{\guillemotleft}}1 {»}{{\guillemotright}}1 {’}{{'}}1}
\lstset{style=emc}

% ---------------------------------------------------------------- titres
\usepackage{titlesec}
\titleformat{\chapter}[display]
  {\sffamily\bfseries}
  {\raggedright\normalsize\color{emcorange}\MakeUppercase{\chaptertitlename}~\thechapter}
  {2pt}
  {\Large\color{emcblue}\raggedright}
  [\vspace{2pt}{\color{emcborder}\titlerule[1.2pt]}]
\titlespacing*{\chapter}{0pt}{-22pt}{10pt}
\titleformat{\section}
  {\sffamily\bfseries\large\color{emcblue}}{\thesection}{0.6em}{}
\titlespacing*{\section}{0pt}{10pt}{4pt}
\titleformat{\subsection}
  {\sffamily\bfseries\normalsize\color{emcslate}}{\thesubsection}{0.6em}{}
\titlespacing*{\subsection}{0pt}{8pt}{2pt}

% ---------------------------------------------------------------- en-têtes
\usepackage{fancyhdr}
\pagestyle{fancy}
\fancyhf{}
\renewcommand{\headrulewidth}{0.4pt}
\renewcommand{\footrulewidth}{0pt}
\fancyhead[L]{\sffamily\scriptsize\color{emcgrey}EMC Helpline — Rapport de stage PFA}
\fancyhead[R]{\sffamily\scriptsize\color{emcgrey}\leftmark}
\fancyfoot[C]{\sffamily\small\color{emcblue}\thepage}
\fancypagestyle{plain}{\fancyhf{}%
  \renewcommand{\headrulewidth}{0pt}%
  \fancyfoot[C]{\sffamily\small\color{emcblue}\thepage}}
\renewcommand{\chaptermark}[1]{\markboth{\thechapter.~#1}{}}

% ---------------------------------------------------------------- glossaire
% \makenoidxglossaries : le glossaire est trié par LaTeX lui-même, sans
% makeglossaries ni Perl — le document compile partout avec xelatex + biber.
\usepackage[acronym,nomain,nopostdot,nogroupskip,nonumberlist,toc=false]{glossaries}
\makenoidxglossaries
\setacronymstyle{long-short}

% ---------------------------------------------------------------- biblio
\usepackage[backend=biber,style=numeric-comp,sorting=none,
            maxbibnames=4,giveninits=true,url=true,doi=false,isbn=false]{biblatex}
\addbibresource{references.bib}
\renewcommand*{\bibfont}{\small}

% ---------------------------------------------------------------- liens
\usepackage[hidelinks,breaklinks=true]{hyperref}
\hypersetup{
  colorlinks=true,linkcolor=emcblue,citecolor=emcorange,urlcolor=emcblue,
  pdftitle={EMC Helpline — Rapport de stage PFA},
  pdfauthor={Sayouba Ouedraogo},
  pdfsubject={Application mobile de signalement des cyberviolences},
  pdfkeywords={Flutter, Firebase, cyberviolence, protection de l'enfance, CMRPI},
  bookmarksnumbered=true,bookmarksopen=true}
\usepackage[french,capitalise,noabbrev]{cleveref}
\usepackage{url}
\urlstyle{sf}

% ---------------------------------------------------------------- utilitaires
\usepackage{xspace}
\usepackage{lastpage}
\newcommand{\emc}{EMC Helpline\xspace}

% Placeholder de capture d'écran.
%   \shot{<largeur>}{<chemin>}{<libellé court>}
% Si le fichier existe, il est inséré ; sinon un cadre pointillé indique
% exactement quel fichier déposer. Le rapport compile donc avant même que la
% moindre capture soit prise.
\newcommand{\shot}[4][2.05]{%
  \IfFileExists{#3}%
    {\includegraphics[width=#2]{#3}}%
    {\begin{tikzpicture}
       \node[draw=emcgrey!60,dashed,line width=0.7pt,fill=emcbg,
             rounded corners=3pt,minimum width=#2,
             minimum height={#1*#2},align=center,inner sep=3pt]
            {\sffamily\scriptsize\color{emcgrey}#4\\[3pt]
             \resizebox{\dimexpr#2-6pt\relax}{!}{\ttfamily\color{emcorange}\detokenize{#3}}};
     \end{tikzpicture}}}

% Capture au format téléphone (ratio 390:844)
\newcommand{\phoneshot}[3]{\shot[2.16]{#1}{#2}{#3}}
% Capture au format écran large (ratio 16:10)
\newcommand{\wideshot}[3]{\shot[0.62]{#1}{#2}{#3}}
